\chapter{Podstawy teoretyczne}

\section{Poker jako gra z niepełną informacją}

Poker jest rodziną gier karcianych, w których wynik zależy zarówno od
losowego rozdania kart, jak i od decyzji podejmowanych przez graczy. W
przeciwieństwie do gier z pełną informacją, takich jak szachy, uczestnicy
rozgrywki nie znają wszystkich elementów stanu gry. Ukryte pozostają przede
wszystkim karty przeciwników, co sprawia, że decyzje podejmowane są w
warunkach niepewności.

Z punktu widzenia sztucznej inteligencji poker jest interesującym problemem
badawczym, ponieważ łączy losowość, niepełną informację, sekwencyjne
podejmowanie decyzji oraz konieczność maksymalizacji wyniku w długim
horyzoncie czasowym. W takim środowisku pojedyncza decyzja nie powinna być
oceniana wyłącznie przez pryzmat wyniku jednego rozdania, lecz przez jej
wpływ na oczekiwaną wartość strategii w wielu powtórzeniach gry.

\section{Texas Hold’em}

Texas Hold’em jest jedną z najpopularniejszych odmian pokera. Każdy gracz
otrzymuje dwie zakryte karty własne (ang. \emph{hole cards}), natomiast na
stole stopniowo odkrywanych jest pięć kart wspólnych (ang. \emph{community
cards}). Gracz tworzy najlepszy możliwy układ pięciokartowy, wykorzystując
dowolną kombinację swoich dwóch kart oraz pięciu kart wspólnych.

Rozgrywka w klasycznym Texas Hold’em składa się z kilku etapów: etapu przed
flopem (ang. \emph{preflop}), flopa (ang. \emph{flop}), czyli odkrycia
trzech pierwszych kart wspólnych, turna (ang. \emph{turn}), czyli czwartej
karty wspólnej, rivera (ang. \emph{river}), czyli piątej i ostatniej karty
wspólnej, oraz showdownu (ang. \emph{showdown}), czyli końcowego porównania
układów. W trakcie kolejnych rund licytacji gracze mogą podejmować decyzje
takie jak czekanie (ang. \emph{check}), postawienie zakładu (ang.
\emph{bet}), sprawdzenie zakładu przeciwnika (ang. \emph{call}), podbicie
(ang. \emph{raise}) lub spasowanie (ang. \emph{fold}). W odmianach no-limit
możliwe jest również stosowanie różnych rozmiarów zakładów, co istotnie
zwiększa przestrzeń możliwych strategii.

\section{Ranking układów pokerowych}

W Texas Hold’em oraz Ultimate Texas Hold’em wynik rozdania zależy od siły
najlepszego układu pięciokartowego utworzonego z dostępnych kart. Układy
pokerowe mają ustaloną hierarchię, która pozwala jednoznacznie porównać rękę
gracza i krupiera podczas showdownu. Przykładowy ranking układów od
najsilniejszego do najsłabszego przedstawiono w
tabeli~\ref{tab:poker_hand_ranking} \cite{tulalip_uth_rules}.

Najsilniejszym układem jest poker królewski, natomiast najsłabszym układem
jest wysoka karta. W przypadku gdy dwóch uczestników rozdania posiada układ
tej samej kategorii, o zwycięstwie decydują wartości kart wchodzących w
skład układu oraz karty dodatkowe, określane jako kickery.

\begin{table}[h]
\centering
\caption{Ranking układów pokerowych od najsilniejszego do najsłabszego}
\label{tab:poker_hand_ranking}
\begin{tabular}{ll}
\toprule
Układ & Opis \\
\midrule
Poker królewski & A, K, Q, J, 10 w tym samym kolorze \\
Poker & Pięć kolejnych kart w tym samym kolorze \\
Kareta & Cztery karty tej samej wartości \\
Full & Trójka oraz para \\
Kolor & Pięć kart w tym samym kolorze \\
Strit & Pięć kolejnych kart \\
Trójka & Trzy karty tej samej wartości \\
Dwie pary & Dwa różne układy par \\
Para & Dwie karty tej samej wartości \\
Wysoka karta & Brak silniejszego układu \\
\bottomrule
\end{tabular}
\end{table}

\section{Ultimate Texas Hold’em}

\subsection{Charakterystyka gry}

Ultimate Texas Hold’em jest kasynową odmianą pokera opartą na zasadach Texas
Hold’em. W przeciwieństwie do klasycznej odmiany wieloosobowej gracz nie
rywalizuje z innymi graczami, lecz z krupierem reprezentującym kasyno.
Zarówno gracz, jak i krupier otrzymują po dwie karty własne, a następnie
korzystają z pięciu kart wspólnych w celu utworzenia najlepszego możliwego
układu pięciokartowego \cite{nevada_uth_rules,tulalip_uth_rules}.

\subsection{Różnice względem klasycznego Texas Hold’em}

Najważniejsza różnica względem klasycznego Texas Hold’em polega na tym, że w
Ultimate Texas Hold’em gracz podejmuje decyzje przeciwko kasynu, a nie
przeciwko innym graczom. Nie występuje więc klasyczna rywalizacja między
wieloma uczestnikami stołu, nie ma konieczności blefowania przeciwników ani
reagowania na ich indywidualne style gry.

Istotną cechą Ultimate Texas Hold’em jest również ograniczona liczba
momentów decyzyjnych. Gracz może wykonać zakład Play (ang. \emph{Play bet})
tylko raz w trakcie rozdania: przed flopem, po flopie albo po odkryciu
wszystkich kart wspólnych. Im wcześniej gracz zdecyduje się na zagranie
agresywne, tym większy zakład może postawić, ale decyzja jest wtedy
podejmowana przy mniejszej liczbie informacji
\cite{wizard_uth,tulalip_uth_rules}.

\begin{table}[h]
\centering
\caption{Porównanie Texas Hold’em i Ultimate Texas Hold’em}
\label{tab:texas_vs_uth}
\begin{tabular}{lll}
\toprule
Cecha & Texas Hold’em & Ultimate Texas Hold’em \\
\midrule
Przeciwnik & inni gracze & krupier / kasyno \\
Liczba graczy & wielu & jeden gracz przeciwko krupierowi \\
Zakłady & zmienne & ustalone mnożniki Ante \\
Blefowanie & istotne & brak klasycznego blefowania \\
Liczba decyzji & duża & ograniczona \\
\bottomrule
\end{tabular}
\end{table}

\subsection{Przebieg rozdania}

Rozdanie rozpoczyna się od postawienia przez gracza obowiązkowych zakładów
Ante (ang. \emph{Ante bet}) oraz Blind (ang. \emph{Blind bet}). Zakłady te
mają zwykle taką samą wartość. Opcjonalnie gracz może również postawić
zakład poboczny (ang. \emph{side bet}), na przykład Trips, który jest
rozliczany niezależnie od wyniku pojedynku z krupierem
\cite{tulalip_uth_rules}.

Po wniesieniu zakładów gracz oraz krupier otrzymują po dwie zakryte karty
własne. Następnie gracz podejmuje pierwszą decyzję. Może zaczekać albo
wykonać zakład Play w wysokości trzykrotności lub czterokrotności zakładu
Ante.

Jeżeli gracz nie wykona zakładu Play przed flopem, odkrywane są trzy
pierwsze karty wspólne, czyli flop. Na tym etapie gracz może ponownie
zaczekać albo wykonać zakład Play w wysokości dwukrotności Ante. Jeżeli
również po flopie gracz zdecyduje się zaczekać, odkrywane są dwie ostatnie
karty wspólne, czyli turn oraz river. Wtedy gracz musi wybrać pomiędzy
spasowaniem a wykonaniem zakładu Play w wysokości jednokrotności Ante
\cite{wizard_uth,tulalip_uth_rules}.

\subsection{Zakłady i decyzje gracza}

W podstawowym modelu gry występują trzy główne rodzaje zakładów: Ante, Blind
oraz Play. Zakłady Ante i Blind są obowiązkowe i stawiane przed rozpoczęciem
rozdania. Zakład Play jest natomiast bezpośrednio związany z decyzją gracza
i może zostać wykonany tylko raz w trakcie rozdania.

Z punktu widzenia modelowania problemu jako środowiska uczenia przez
wzmacnianie najważniejszy jest zakład Play, ponieważ jego wysokość i moment
wykonania wynikają z decyzji agenta. Ante i Blind można traktować jako
początkowy koszt udziału w rozdaniu, natomiast Play odzwierciedla wybór
strategii w aktualnym stanie gry.

Dostępne decyzje gracza w poszczególnych etapach rozdania przedstawiono w
tabeli~\ref{tab:uth_decisions}.

\begin{table}[h]
\centering
\caption{Dostępne decyzje gracza w Ultimate Texas Hold’em}
\label{tab:uth_decisions}
\begin{tabular}{lll}
\toprule
Etap & Dostępne decyzje & Wysokość zakładu Play \\
\midrule
Preflop & czekanie lub zakład Play & 3x albo 4x Ante \\
Flop & czekanie lub zakład Play & 2x Ante \\
Po odkryciu całego stołu & spasowanie lub zakład Play & 1x Ante \\
\bottomrule
\end{tabular}
\end{table}

\subsection{Rozstrzyganie rozdania}

Po zakończeniu etapu decyzyjnego krupier odsłania swoje dwie karty własne.
Zarówno gracz, jak i krupier tworzą najlepszy możliwy układ pięciokartowy z
siedmiu dostępnych kart: dwóch kart własnych oraz pięciu kart wspólnych.
Następnie układy są porównywane zgodnie ze standardowym rankingiem układów
pokerowych.

Krupier kwalifikuje się do gry, jeżeli posiada co najmniej parę. Jeżeli
krupier się nie kwalifikuje, zakład Ante jest zwracany graczowi. Zakład Play
jest rozliczany na podstawie bezpośredniego porównania układu gracza i
krupiera. W przypadku zwycięstwa gracza Play wypłacany jest w stosunku 1:1,
w przypadku przegranej gracz traci zakład Play, a w przypadku remisu zakład
jest zwracany \cite{tulalip_uth_rules,wizard_uth}.

Zakład Blind ma odrębną logikę rozliczania. Jeżeli gracz pokona krupiera
układem o wartości co najmniej strita, Blind wypłacany jest zgodnie z tabelą
wypłat. Jeżeli gracz wygra rozdanie układem słabszym niż strit, zakład Blind
jest zwracany. W przypadku przegranej gracza zakład Blind jest tracony, a w
przypadku remisu zwracany \cite{tulalip_uth_rules}.

\subsection{Zakłady poboczne}

W wielu wariantach Ultimate Texas Hold’em dostępne są również zakłady
poboczne, takie jak Trips lub Progressive. Zakład Trips jest opcjonalny i
wypłacany na podstawie końcowego układu gracza, niezależnie od tego, czy
gracz pokonał krupiera. Najczęściej wygrywa on wtedy, gdy końcowy układ
gracza ma wartość co najmniej trójki \cite{tulalip_uth_rules}.

Zakłady poboczne mogą różnić się w zależności od kasyna i stosowanej tabeli
wypłat.

% NOTE: Trips i Progressive pominięto zgodnie z zakresem środowiska.
W zaimplementowanym środowisku nie uwzględniono zakładów pobocznych Trips
ani Progressive. Wynika to z faktu, że głównym celem pracy jest modelowanie
sekwencyjnego procesu decyzyjnego dotyczącego zakładu Play. Zakłady poboczne
są opcjonalne, zależą od tabel wypłat konkretnego kasyna i zwiększają
wariancję wyniku, nie zmieniając podstawowej struktury decyzji: czekanie,
zakład Play albo spasowanie.

\subsection{Tabele wypłat}

Zakład Blind nie jest wypłacany dla każdego zwycięskiego układu gracza.
Zgodnie z tabelą wypłat, bonus na zakładzie Blind wypłacany jest wtedy, gdy
gracz pokona krupiera i posiada układ o wartości co najmniej strita. Jeżeli
gracz pokona krupiera układem słabszym niż strit, zakład Blind jest
zwracany. Tabelę wypłat zakładu Blind przedstawiono w
tabeli~\ref{tab:blind_payouts} \cite{tulalip_uth_rules}.

\begin{table}[h]
\centering
\caption{Tabela wypłat zakładu Blind przyjęta w modelu środowiska}
\label{tab:blind_payouts}
\begin{tabular}{ll}
\toprule
Układ gracza & Wypłata \\
\midrule
Poker królewski & 500:1 \\
Poker & 50:1 \\
Kareta & 10:1 \\
Full & 3:1 \\
Kolor & 3:2 \\
Strit & 1:1 \\
Pozostałe układy & Zwrot \\
\bottomrule
\end{tabular}
\end{table}

Zakład Trips oraz zakład Progressive są opcjonalnymi zakładami pobocznymi,
rozliczanymi niezależnie od pojedynku z krupierem i zależnymi od tabel
wypłat konkretnego kasyna. Ponieważ nie wpływają one na podstawową decyzję
Play, w zaimplementowanym środowisku zakładów tych nie uwzględniono.

\section{Wartość oczekiwana i przewaga kasyna}

Wartość oczekiwana (ang. \emph{expected value}, EV) określa średni wynik
finansowy decyzji lub strategii przy założeniu dużej liczby powtórzeń gry. W
kontekście Ultimate Texas Hold’em pojedyncze rozdanie może zakończyć się
zarówno zyskiem, jak i stratą, jednak jakość strategii należy oceniać na
podstawie średniego wyniku uzyskiwanego w wielu rozdaniach.

Przewaga kasyna (ang. \emph{house edge}) jest matematyczną miarą określającą
oczekiwany zysk kasyna względem zakładu gracza. Nie oznacza ona, że kasyno
wygrywa każde pojedyncze rozdanie, lecz że przy dużej liczbie rozegranych
rozdań średni wynik będzie przesuwał się na korzyść kasyna. Gracz może więc
osiągać krótkoterminowe zyski, jednak w długim okresie wynik gry powinien
zbliżać się do wartości oczekiwanej wynikającej z zasad gry
\cite{investopedia_house_edge}.

W Ultimate Texas Hold’em przewaga kasyna jest zwykle podawana względem
początkowego zakładu Ante. Ponieważ gracz może dodatkowo zwiększyć swoje
zaangażowanie poprzez zakład Play, obok klasycznej przewagi kasyna stosuje
się również miarę określaną jako element ryzyka (ang. \emph{element of
risk}), która odnosi oczekiwaną stratę do średniej całkowitej kwoty
zaangażowanej w rozdanie \cite{wizard_house_edge,wizard_uth}. Nawet
niewielka procentowa przewaga prowadzi więc do istotnej oczekiwanej straty w
długim horyzoncie czasowym.

Przykładowe wartości przewagi kasyna dla wybranych gier przedstawiono w
tabeli~\ref{tab:house_edge_examples} \cite{wizard_house_edge}.

\begin{table}[h]
\centering
\caption{Przykładowe wartości przewagi kasyna dla wybranych gier kasynowych}
\label{tab:house_edge_examples}
\begin{tabular}{ll}
\toprule
Gra & Przewaga kasyna \\
\midrule
Blackjack & 0{,}28\% \\
Ultimate Texas Hold’em & 2{,}19\% \\
Ruletka europejska & 2{,}70\% \\
Ruletka amerykańska & 5{,}26\% \\
Slot machines & 2--15\% \\
\bottomrule
\end{tabular}
\end{table}

Tabela~\ref{tab:house_edge_examples} pokazuje, że przewaga kasyna jest cechą
konstrukcyjną gier hazardowych, ale jej poziom istotnie różni się pomiędzy
gramim, a podane wartości zależą przy tym od wariantu zasad gry oraz przyjętego
sposobu normalizacji. Celem niniejszej pracy nie jest wykazanie, że agent
może trwale uzyskać dodatnią wartość oczekiwaną w grze zaprojektowanej na
korzyść kasyna, lecz sprawdzenie, które metody uczenia przez wzmacnianie są
w stanie nauczyć się strategii minimalizującej tę przewagę.

\section{Strategie bazowe i ocena jakości strategii}

Strategie bazowe stanowią punkty odniesienia służące do oceny agentów
uczących się. Agent losowy wybiera jedną z legalnych akcji bez uwzględniania
siły kart i pozwala zweryfikować poprawność infrastruktury eksperymentalnej.
Agent regułowy wykorzystuje z góry określone heurystyki i reprezentuje
bardziej wymagający punkt odniesienia.

Wyniki agentów uczących się powinny być porównywane zarówno pomiędzy sobą,
jak i względem strategii bazowych. Samo przewyższenie agenta losowego nie
jest wystarczającym dowodem uzyskania jakościowej strategii, natomiast
porównanie z agentem regułowym pozwala ocenić, czy uczenie prowadzi do
decyzji konkurencyjnych wobec jawnie zaprojektowanych zasad.

\section{Uczenie przez wzmacnianie}

Uczenie przez wzmacnianie (ang. \emph{reinforcement learning}, RL) jest
paradygmatem uczenia maszynowego, w którym agent uczy się podejmowania
decyzji poprzez interakcję ze środowiskiem. W każdym kroku agent obserwuje
aktualny stan (ang. \emph{state}) środowiska, wybiera akcję (ang.
\emph{action}), a następnie otrzymuje nagrodę (ang. \emph{reward}) oraz
przechodzi do kolejnego stanu. Celem agenta jest wyuczenie strategii,
określanej również jako polityka decyzyjna (ang. \emph{policy}), która
maksymalizuje sumę nagród w długim horyzoncie czasowym
\cite{sutton_barto_rl}.

Z perspektywy uczenia przez wzmacnianie Ultimate Texas Hold’em jest
interesującym środowiskiem, ponieważ decyzje gracza są sekwencyjne i
podejmowane przy różnym poziomie dostępnej informacji. Przed flopem gracz
zna jedynie swoje dwie karty, ale może wykonać najwyższy zakład Play. Po
flopie posiada więcej informacji, lecz maksymalny zakład jest mniejszy. Po
odkryciu wszystkich kart wspólnych zna już pełny board, ale może wykonać
wyłącznie zakład 1x Ante albo spasować. Agent musi więc nauczyć się
kompromisu pomiędzy ryzykiem, ilością informacji oraz potencjalną wypłatą.

\subsection{Podstawowe pojęcia}

Zachowanie agenta opisuje polityka \(\pi\), która może być deterministyczna,
przypisując stanowi jedną akcję, albo stochastyczna, określając rozkład
prawdopodobieństwa wyboru akcji. Jakość decyzji ocenia się nie na podstawie
pojedynczej nagrody, lecz zwrotu epizodycznego, czyli sumy nagród uzyskanych
do końca epizodu:

\begin{equation}
    G_t = \sum_{k=0}^{T-t-1} \gamma^{k}\, r_{t+k},
\end{equation}

gdzie \(\gamma \in [0,1]\) jest współczynnikiem dyskontowania, a \(T\)
oznacza koniec epizodu. Współczynnik \(\gamma\) reguluje wagę nagród
odległych w czasie względem nagród natychmiastowych \cite{sutton_barto_rl}.

Oczekiwany zwrot z danego stanu opisuje funkcja wartości stanu
\(V^{\pi}(s)\), natomiast oczekiwany zwrot po wykonaniu akcji w danym stanie
- funkcja wartości akcji \(Q^{\pi}(s,a)\). Znajomość tych funkcji pozwala
porównywać decyzje i stanowi podstawę wielu metod uczenia
\cite{sutton_barto_rl}.

W trakcie uczenia agent musi równoważyć eksplorację, czyli wypróbowywanie
nowych akcji w celu zdobycia informacji, oraz eksploatację, czyli wybór
akcji uznanych dotąd za najlepsze. Niewłaściwy kompromis między nimi
prowadzi do utknięcia w strategii suboptymalnej albo do niestabilnego
przebiegu uczenia.

Funkcje wartości i polityki można reprezentować w postaci tablicowej, gdy
przestrzeń stanów jest niewielka, albo za pomocą aproksymacji funkcji, na
przykład sieci neuronowych, gdy przestrzeń stanów jest duża lub ciągła.
Aproksymacja umożliwia generalizację pomiędzy podobnymi stanami, lecz może
pogarszać stabilność i zbieżność uczenia \cite{sutton_barto_rl}.

W rozważanym problemie szczególne znaczenie mają dwa elementy konfiguracji
procesu uczenia: sposób reprezentacji obserwacji, decydujący o tym, jakie
informacje trafiają na wejście agenta, oraz konstrukcja funkcji nagrody,
przekształcająca finansowy wynik rozdania w sygnał uczący. Oba te elementy
są przedmiotem analizy w dalszej części pracy.

\section{Modelowanie UTH jako procesu decyzyjnego z częściową obserwowalnością}

W Ultimate Texas Hold’em agent nie obserwuje pełnego stanu środowiska,
ponieważ nie zna kart krupiera, kart spalonych ani przyszłej kolejności
talii. Podejmuje decyzję na podstawie obserwacji zawierającej własne karty,
odkryte karty wspólne, bieżący etap rozdania oraz legalne akcje.

Problem można zatem traktować jako częściowo obserwowalny proces decyzyjny.
Polityka agenta odwzorowuje dostępną obserwację na rozkład
prawdopodobieństwa wyboru akcji, natomiast celem uczenia jest maksymalizacja
oczekiwanego zwrotu uzyskiwanego w wielu epizodach. Rozróżnienie pełnego
stanu środowiska od obserwacji agenta jest kluczowe dla zachowania niepełnej
informacji charakterystycznej dla pokera i zostało bezpośrednio odwzorowane
w projekcie środowiska opisanym w kolejnym rozdziale.